% The clock tree at the instant of wake, beside the one the application expects.
\begin{tikzpicture}[font=\sffamily\small]
% ================= left: the instant after wake
\begin{scope}
\node[font=\sffamily\bfseries\small, anchor=west] at (-0.4,2.1) {at the instant of wake};
\node[osc, text width=20mm, fill=black!8] (hse1) at (0,0)     {HSE 8 MHz\\{\tiny off}};
\node[osc, text width=20mm, fill=hwcol]   (hsi1) at (0,-1.1)  {HSI\\{\tiny selected}};
\node[pll, text width=20mm, fill=black!8] (pll1) at (2.6,0)    {PLL1\\{\tiny off}};
\node[clkmux] (mux1) at (2.6,-1.1) {sel};
\node[clkdom, text width=22mm, fill=black!8] (sys1) at (5.2,-1.1) {SYSCLK\\{\tiny reduced speed}};
\draw[clkline] (hsi1) -- (mux1);
\draw[clkline, draw=black!35] (hse1) -- (pll1);
\draw[clkline, draw=black!35] (pll1) -- (mux1);
\draw[clkline] (mux1) -- (sys1);
\node[clkdom, text width=22mm, fill=hwcol] (vos1) at (5.2,0.1) {voltage scale\\{\tiny the stop scale}};
\node[clkdom, text width=22mm, fill=hwcol] (lat1) at (5.2,1.05) {flash latency\\{\tiny for the low scale}};
\end{scope}

% ================= divider
\draw[dashed, draw=linecol] (7.2,-2.6) -- (7.2,2.4);

% ================= right: what the application expects
\begin{scope}[xshift=7.8cm]
\node[font=\sffamily\bfseries\small, anchor=west] at (-0.4,2.1) {what the application expects};
\node[osc, text width=20mm, fill=hwcol] (hse2) at (0,0)    {HSE 8 MHz\\{\tiny bypass, on}};
\node[osc, text width=20mm, fill=black!8] (hsi2) at (0,-1.1) {HSI\\{\tiny idle}};
\node[pll, text width=20mm] (pll2) at (2.6,0) {PLL1\\{\tiny $\times$140, /2, /2}};
\node[clkmux] (mux2) at (2.6,-1.1) {sel};
\node[clkdom, text width=22mm, fill=usrcol] (sys2) at (5.2,-1.1) {SYSCLK\\{\tiny 280 MHz}};
\draw[clkline] (hse2) -- (pll2);
\draw[clkline] (pll2) -- (mux2);
\draw[clkline, draw=black!35] (hsi2) -- (mux2);
\draw[clkline] (mux2) -- (sys2);
\node[clkdom, text width=22mm, fill=usrcol] (vos2) at (5.2,0.1) {voltage scale\\{\tiny the run scale}};
\node[clkdom, text width=22mm, fill=usrcol] (lat2) at (5.2,1.05) {flash latency\\{\tiny for 280 MHz}};
\end{scope}

% ================= the restore, drawn as the ordered path between them
\draw[flow, draw=red!70!black] (4.6,-2.05) -- node[lbl, above]
  {restore\_after\_stop()} (12.3,-2.05);
\node[umlnote, text width=62mm, anchor=north west] at (0.6,-2.9)
  {Nothing in the silicon performs that transition. The order inside it is voltage scale, then flash wait states, then the phase-locked loop, then the switch. Reversing it is the classic way to produce a board that runs at reset speed or not at all.};
\end{tikzpicture}
